% Generated by scripts/assemble.py; corrections belong in source/corrections.json.
\examhead{51}{上海财经大学 2025 年数学分析试卷}{52}
% p053-b12; source PDF page 53
\par\Needspace{5\baselineskip}
\blockanchor{p053-b12}
\majorquestion{一}计算题.\par
% p053-b13; source PDF page 53
\blockanchor{p053-b13}
\question{1}(14 分) 计算 $\lim_{n\to \infty}(n + 1)! \left(e - 1 - \frac{1}{1!} - \frac{1}{2!} - \cdots - \frac{1}{n!}\right)$.\par
% p053-b14; source PDF page 53
\blockanchor{p053-b14}
\question{2}(14 分) 计算  $\lim_{n\to\infty}n\left(\frac{\sin\frac{\pi}{n}}{n^{2}+1}+\frac{\sin\frac{2\pi}{n}}{n^{2}+2}+\cdots+\frac{\sin\frac{n\pi}{n}}{n^{2}+n}\right)$ .\par
% p053-b15; source PDF page 53
\blockanchor{p053-b15}
\question{3}(14 分) 求 $\sum_{n=1}^{\infty} n(n+1)x^{n}$ 的和函数.\par
% p053-b16; source PDF page 53
\blockanchor{p053-b16}
\question{4}(14 分) 计算 $\iint_{D} \sqrt{\frac{1 - x^2 - y^2}{1 + x^2 + y^2}} \, \mathrm{d}x \, \mathrm{d}y$, 其中 $D = \{(x, y) \mid x^2 + y^2 \leq 1\}$ 在第一象限的部分.\par
% p053-b17; source PDF page 53
\blockanchor{p053-b17}
\question{5}(14 分) 计算 $\int_{L} (xy + yz) \, \mathrm{d}s$, 其中 $L$ 是球面 $x^{2} + y^{2} + z^{2} = 1$ 与平面 $x + y + z = 0$ 的交线.\par
% p053-b18; source PDF page 53
\par\Needspace{5\baselineskip}
\blockanchor{p053-b18}
\majorquestion{二}解答与证明.\par
% p053-b19; source PDF page 53
\blockanchor{p053-b19}
\question{6}(10 分) 已知 $\lim_{n\to \infty}(a_{n + 2} - a_n) = a,$ 问 $\lim_{n\to \infty}\frac{a_n}{n}$ 是否存在, 若存在, 请给出证明以及它的值.\par
% p053-b20; source PDF page 53
\blockanchor{p053-b20}
\question{7}(10 分) 设 $f(x)$ 在 $[0,3]$ 上可导, $f(0) + f(1) + f(2) = 3$, $f(3) = 1$. 求证: 存在 $\xi \in (0,3)$, 使得 $f'(\xi) = 0$.\par
% p053-b21; source PDF page 53
\blockanchor{p053-b21}
\question{8}(10 分) 函数列 $\{f_n(x)\}$ 在 $[a, b]$ 上连续, 在 $(a, b)$ 上一致收敛. 证明: $\{f_n(x)\}$ 在点 $a$ 和 $b$ 处收敛, 并说明极限函数在区间 $[a, b]$ 上连续.\par
% p053-b22; source PDF page 53
\blockanchor{p053-b22}
\question{9}(10 分) 已知 $\sum_{n=1}^{\infty} na_n$ 收敛, 判断 $\sum_{n=1}^{\infty} (n+100)a_n$ 的敛散性.\par
% p053-b23; source PDF page 53
\par\Needspace{6\baselineskip}
\blockanchor{p053-b23}
\question{10}(10 分) 设 $f(x)$ 在 $[0,1]$ 上连续且单调递增, $g(x)$ 在 $[0,1]$ 上连续且单调递减. 证明:\par
\nopagebreak[4]
% p053-b24; source PDF page 53
\blockanchor{p053-b24}
\displayformula{\int_ {0} ^ {1} f (x) g (x) \mathrm{d} x \leq \int_ {0} ^ {1} f (x) \mathrm{d} x \int_ {0} ^ {1} g (x) \mathrm{d} x.}
% p054-b00; source PDF page 54
\par\Needspace{6\baselineskip}
\blockanchor{p054-b00}
\question{11}(10 分) 设\par
\nopagebreak[4]
% p054-b01; source PDF page 54
\blockanchor{p054-b01}
\displayformula{f (x, y) = \left\{ \begin{array}{l l} (x ^ {2} + y ^ {2}) \sin \frac {1}{x ^ {2} + y ^ {2}}, & (x, y) \neq (0, 0); \\ 0, & (x, y) = (0, 0). \end{array} \right.}
% p054-b02; source PDF page 54
\blockanchor{p054-b02}
判断  $f(x,y)$  在  $(0,0)$  处是否可微，偏导数是否连续.\par
% p054-b03; source PDF page 54
\par\Needspace{6\baselineskip}
\blockanchor{p054-b03}
\question{12}(10 分) 设 $\Sigma$ 是单位球面 $x^{2} + y^{2} + z^{2} = 1$, $f$ 是一个连续函数, 证明:\par
\nopagebreak[4]
% p054-b04; source PDF page 54
\blockanchor{p054-b04}
\displayformula{\iint_ {\Sigma} f (x + y + z) \mathrm{d} S = 2 \pi \int_ {- 1} ^ {1} f (\sqrt {3} u) \mathrm{d} u.}
% p054-b05; source PDF page 54
\blockanchor{p054-b05}
\question{13}(10 分) 计算 $f(x, y, z) = \ln x + 2\ln y + 4\ln z$ 在 $x^2 + y^2 + z^2 = 7R^2 (x, y, z > 0, R > 0)$ 下的最大值, 并证明对任意的正数 $a, b, c$, 都有 $ab^2c^4 \leq 1024\left(\frac{a + b + c}{7}\right)^7$.\par
